%% kapa: exam=1 round=29 year=2018 period=1 form=A subject=economics subject_ko=경제학원론 seq=31 number=71
%% source: Q-Net 감정평가사 기출문제 (cst00309) artlSeq=5202808 file=extracted/1008/5202808/2018년 제29회 감정평가사 1차시험 기출문제/2018년 제29회 감정평가사 1차시험 기출문제/1교시 A형.pdf
%% conversion: kordoc 4.24.51 -> tools/gen1.py (auto); review=proofread-sonnet; flags=none
\begin{question}{71}
甲국의 국민소득( $Y$)은 소비( $C$), 민간투자( $I$), 정부지출( $G$), 순수출( $NX$)의 합과 같다. 2016년과 같이 2017년에도 조세( $T$)와 정부지출의 차이( $T-G$)는 음(-)이었고 절대크기는 감소하였으며, 순수출은 양(+)이었지만 절대크기는 감소하였다. 이로부터 유추할 수 있는 2017년의 상황으로 옳은 것을 모두 고른 것은?
\begin{framedtext}
ㄱ. 국가채무는 2016년 말에 비해 감소하였다. \par
ㄴ. 순대외채권은 2016년 말에 비해 감소하였다. \par
ㄷ. 민간저축은 민간투자보다 더 많았다. \par
ㄹ. 민간저축과 민간투자의 차이는 2016년보다 그 절대크기가 감소하였다. \par
\end{framedtext}
\begin{choices}
  \item ㄱ, ㄴ
  \item ㄱ, ㄷ
  \item ㄴ, ㄷ
  \item ㄴ, ㄹ
  \item ㄷ, ㄹ
\end{choices}
\end{question}
\answer{5}
